\documentclass[exercice,a4]{thedocument}%
\usepackage[french]{babel}%
\usepackage{tabularx}%
\usepackage{array}%
\startdatakeys
    \source{22GENMATNC1}
    \BO{2026}
    \cycle{4}
    \competencesDuSocle{CA,RA}
    \objectifsApprentissage{C4C2ab,C4D2ae}
    \motsCles{fréquence,pourcentage,médiane}
\enddatakeys
\begin{document}%


\medskip

On a relevé la vitesse du vent à 13 heures du 1\up{er} au 15 novembre sur une plage de Nouvelle- Calédonie.

Les vitesses approchées sont données, en nœuds, dans le tableau ci-dessous :

\begin{center}
\begin{tabularx}{\linewidth}{|m{3cm}|*{15}{>{\centering \arraybackslash}X|}}\hline
Jours du 1\up{er} au 15 novembre&1 &2 &3 &4 &5 &6 &7 &8 &9 &10 &11 &12 &13 &14 &15\\ \hline
Vitesse du vent en nœuds& 10 &15 &20 &20 &15 &10 &10 &20 &15 &25 &25 &25 &20 &15 &15\\ \hline
\end{tabularx}
\end{center}

\begin{enumerate}
\item À partir des données ci-dessus, compléter le tableau figurant sur l'annexe.
\item Calculer le pourcentage de jours où la vitesse de vent est supérieure ou égale à 15 nœuds sur la plage, entre le 1\up{er} et le 15 novembre.
\item Déterminer la vitesse médiane du vent sur la plage durant cette période.
\end{enumerate}

\bigskip


\begin{center}\textbf{ANNEXE À RENDRE AVEC LA COPIE} \end{center}

\medskip

\textbf{Exercice 3 : question 1}

\begin{center}
\begin{tabularx}{\linewidth}{|m{6.5cm}|*{4}{>{\centering \arraybackslash}X|}}\hline
Vitesse du vent (en nœuds)	&10 &15 &20 &25\\ \hline
Nombre de jours				&3	&	&	&3\\ \hline
Fréquence en \,\% arrondie à l'unité&&33&&\\ \hline
\end{tabularx}
\end{center}

\bigskip

\showthesource
\end{document}
